% ============================================================
% EXPANDED LITERATURE REVIEW — 12 Key Papers + Review
% For defense_presentation.tex (Beamer)
% ============================================================

% ============================================================
% SLIDE 6 (REPLACEMENT) — Literature Review: 12 Key Works
% ============================================================
\begin{frame}{Literature Review: 12 Foundational Works}
  \centering
  \footnotesize
  \begin{tabular}{p{0.8cm}p{3.8cm}p{5.8cm}}
    \toprule
    \textbf{Year} & \textbf{Authors} & \textbf{Key Contribution \& Relevance} \\
    \midrule
    1989 & McCloskey \& Cohen & Catastrophic forgetting in neural networks — core motivation for our 15\% linguistic regularisation \\
    1997 & Caruana & Multi-task learning foundations — basis of our hybrid NER + linguistic training \\
    2001 & Lafferty et al. & Conditional Random Fields (CRF) — first statistical NER standard \\
    2015 & Huang et al. & BiLSTM-CRF — neural sequence labeling breakthrough \\
    2018 & Hellwig \& Nehrdich & Character-level Sanskrit segmentation — early deep learning for Sanskrit \\
    2019 & Devlin et al. & BERT — Transformer pretraining revolution \\
    2022 & Xue et al. & ByT5 — byte-level tokenization (our base architecture choice) \\
    2023 & Sarkar et al. & First CRF-based Sanskrit NER attempt — showed data scarcity problem \\
    2023 & Dettmers et al. & QLoRA — enabled our single-GPU 594M model training \\
    2024 & Nehrdich et al. & ByT5-Sanskrit — first large-scale Sanskrit language model (our base) \\
    2025 & Sarkar et al. & Mahānāma — first large gold-standard Sanskrit NER corpus (our D1 benchmark) \\
    2025 & Hellwig \& Biagetti & Sanskrit Sembank — source of our novel silver NER extraction \\
    \bottomrule
  \end{tabular}
  \note{These 12 works form the foundation of our methodology. We combine ideas from statistical NER, modern efficient fine-tuning (QLoRA), and Sanskrit-specific resources (Mahānāma + Sembank).}
\end{frame}

% ============================================================
% SLIDE 7 — Research Gap (Updated)
% ============================================================
\begin{frame}{Research Gap: What Was Missing}
  \begin{columns}[T]
    \begin{column}{0.52\textwidth}
      \textbf{What existed before this thesis:}
      \begin{itemize}
        \item Strong general NER methods (CRF, BERT, ByT5)
        \item Good Sanskrit resources (DCS, Mahānāma 2025)
        \item Efficient fine-tuning (LoRA/QLoRA)
      \end{itemize}
      \vspace{0.3cm}
      \textbf{What was missing:}
      \begin{itemize}
        \item No model that satisfies \textbf{all three desiderata} at once
        \item No correct way to extract NER labels from Sembank
        \item No evaluation that checks linguistic preservation + cross-domain
      \end{itemize}
    \end{column}
    \begin{column}{0.44\textwidth}
      \centering
      \begin{tikzpicture}[
        box/.style={draw=navy, fill=navy!8, thick, rounded corners, minimum width=3.2cm, minimum height=1.1cm, align=center, font=\footnotesize},
        gap/.style={draw=alertred, fill=alertred!15, thick, rounded corners, minimum width=3.2cm, minimum height=1.1cm, align=center, font=\footnotesize\bfseries}
      ]
        \node[box] (d1) at (0,2.2) {High NER Accuracy};
        \node[box] (d2) at (0,0.9) {Linguistic Capability};
        \node[box] (d3) at (0,-0.4) {Cross-Domain};
        \node[gap] at (0,-1.7) {All Three Together};
        \draw[->, thick, gold] (d1) -- (d2) -- (d3);
      \end{tikzpicture}
    \end{column}
  \end{columns}

  \vspace{0.2cm}
  \begin{alertblock}{Our Contribution}
    We are the \textbf{first} to combine all three using 15\% linguistic regularisation + corrected Sembank silver data + three-dimensional evaluation.
  \end{alertblock}
  \note{This is the precise gap our thesis fills. Previous work always sacrificed at least one of the three requirements.}
\end{frame}

% ============================================================
% OPTIONAL: SHORT REVIEW PARAGRAPH (can be used as speaker note or extra slide)
% ============================================================
\begin{frame}{Literature Review: Synthesis}
  \textbf{Synthesis of 12 Works:}

  The field evolved from statistical models (CRF, 2001) to neural sequence taggers (BiLSTM-CRF, 2015) to large pretrained Transformers (BERT 2019, ByT5 2022). Sanskrit-specific work remained limited until 2023–2025, when Mahānāma and ByT5-Sanskrit finally provided the necessary scale.

  Our work sits at the intersection of three lines:
  \begin{enumerate}
    \item \textbf{Efficient adaptation} (QLoRA 2023)
    \item \textbf{Sanskrit resources} (Mahānāma + Sembank 2025)
    \item \textbf{Multi-task regularisation} (Caruana 1997 + catastrophic forgetting literature)
  \end{enumerate}

  This combination had not been attempted before for Sanskrit NER.
  \note{This short review can be used if the examiner asks for a broader perspective on how our work relates to the literature.}
\end{frame}

% ============================================================
% END
% ============================================================